\renewcommand{\arraystretch}{1.05}
\footnotesize
\onecolumn
\begin{longtable}{@{}r p{9cm} c c c c@{}}
\caption{Full 673-record title/abstract screening validation sample (full census of the 276 pipeline-INCLUDE records plus a 397-record probability sample of the pipeline-EXCLUDE stratum, exceeding the Cochran/FPC minimum; pipeline-UNCERTAIN/ERROR records are out of scope; see Section~\ref{sec:validation}). \textbf{Rev.\ A} / \textbf{Rev.\ B}: independent human-annotator decisions (I = Include, E = Exclude; blank where annotation is still pending). \textbf{Final}: adjudicated label used as ground truth (marked $^{\dagger}$ where adjudication was required to resolve a Rev.\ A/Rev.\ B disagreement). \textbf{Pipeline}: the original stage-1 screening pipeline decision from \texttt{results/stage1\_results.csv}. As of this draft, 0 of 673 records have a resolved Final label.}
\label{tab:screening_validation_full} \\
\toprule
\textbf{\#} & \textbf{Title} & \textbf{Rev.\,A} & \textbf{Rev.\,B} & \textbf{Final} & \textbf{Pipeline} \\
\midrule
\endfirsthead
\multicolumn{6}{c}{\tablename~\thetable{} -- \textit{continued from previous page}} \\
\toprule
\textbf{\#} & \textbf{Title} & \textbf{Rev.\,A} & \textbf{Rev.\,B} & \textbf{Final} & \textbf{Pipeline} \\
\midrule
\endhead
\midrule
\multicolumn{6}{r}{\textit{Continued on next page}} \\
\endfoot
\bottomrule
\endlastfoot
1 & Subtraction-free artifact-aware digital subtraction angiography image generation for head and neck vessels from motion data. &  &  &  & I \\
2 & Zero-fluoroscopy approach for ablation of supraventricular tachycardia using the Ensite NavX system: a multicenter experience &  &  &  & E \\
3 & Emerging Image-Guided Navigation Techniques for Cardiovascular Interventions: A Scoping Review. &  &  &  & E \\
4 & Intergrader Agreement on Qualitative and Quantitative Assessment of Diabetic Retinopathy Severity Using Ultra-Widefield Imaging: INSPIRED Study Report 1. &  &  &  & E \\
5 & Diagnosis of Polypoidal Choroidal Vasculopathy From Fluorescein Angiography Using Deep Learning. &  &  &  & I \\
6 & Hyperparameter optimization of YOLO models for invasive coronary angiography lesion detection and assessment. &  &  &  & I \\
7 & Application of Pseudo-Three-Dimensional Residual Network to Classify the Stages of Moyamoya Disease. &  &  &  & I \\
8 & Optical coherence tomography angiography analysis of interocular asymmetry in hydroxychloroquine retinal toxicity. &  &  &  & E \\
9 & End-to-end deep learning approach for uterine artery-ovarian artery anastomosis detection from digital subtraction angiography sequences. &  &  &  & I \\
10 & Deep Learning-Based Intraoperative Stent Graft Segmentation on Completion Digital Subtraction Angiography During Endovascular Aneurysm Repair. &  &  &  & I \\
11 & Clinical application of automated optical coherence tomography angiography for retinal fluid segmentation: a study of real-world data. &  &  &  & E \\
12 & A Deep Learning Algorithm for Classifying Diabetic Retinopathy Using Optical Coherence Tomography Angiography. &  &  &  & E \\
13 & Semi-supervised segmentation of coronary DSA using mixed networks and multi-strategies. &  &  &  & I \\
14 & A dual-task deep learning framework for automated detection and classification of coronary artery lesions in invasive coronary angiography imaging. &  &  &  & I \\
15 & Validation of UniverSeg for Interventional Abdominal Angiographic Segmentation. &  &  &  & I \\
16 & Fully automated Bayesian analysis for quantifying the extent and distribution of pulmonary perfusion changes on CT pulmonary angiography in CTEPH. &  &  &  & E \\
17 & Pericoronary Fat Attenuation: Diagnosis and Clinical Implications. &  &  &  & E \\
18 & Development and Validation of a Deep Learning-Based Segmentation Method for Fenestration Marker and Graft Body Identification in Fenestrated Endovascular Aortic Repair. &  &  &  & I \\
19 & Detecting Large Vessel Occlusion at Multiphase CT Angiography by Using a Deep Convolutional Neural Network. &  &  &  & E \\
20 & Surgical navigation for guidewire placement from intraoperative fluoroscopy in orthopaedic surgery. &  &  &  & I \\
21 & FluoroSAM: A Language-promptable Foundation Model for Flexible X-ray Image Segmentation. &  &  &  & I \\
22 & Development of Novel Artificial Intelligence to Detect the Presence of Clinically Meaningful Coronary Atherosclerotic Stenosis in Major Branch from Coronary Angiography Video. &  &  &  & I \\
23 & Enhancing accuracy of detecting left atrial dilatation on CT pulmonary angiography. &  &  &  & E \\
24 & Diagnostic performance of deep learning algorithm for analysis of computed tomography myocardial perfusion. &  &  &  & E \\
25 & Stenosis degree and plaque burden differ between the major epicardial coronary arteries supplying ischemic territories. &  &  &  & E \\
26 & Exploration of oral microbiota alteration and AI-driven non-invasive hyperspectral imaging for CAD prediction. &  &  &  & E \\
27 & An artificial intelligence-based model for exploration of coronary artery dissection and planning of percutaneous coronary intervention &  &  &  & I \\
28 & X-Ray to CT Rigid Registration Using Scene Coordinate Regression &  &  &  & I \\
29 & Unsupervised machine learning for risk stratification and identification of relevant subgroups of ascending aorta dimensions using cardiac CT and clinical data. &  &  &  & E \\
30 & A Temporary Transformer Network for Guide- Wire Segmentation &  &  &  & I \\
31 & Deep learning-driven automated quality assessment of ultra-widefield optical coherence tomography angiography images for diabetic retinopathy &  &  &  & E \\
32 & Analysis of Video Retinal Angiography With Deep Learning and Eulerian Magnification &  &  &  & I \\
33 & Deep learning-based assessment of coronary artery disease using curved-multiplanar reconstruction: comprehensive evaluation of stenosis, calcification, and plaque. &  &  &  & E \\
34 & Radiation Safety in Musculoskeletal Radiology: Addressing Challenges and Implementing Solutions &  &  &  & E \\
35 & Automated Prediction of Proximal Middle Cerebral Artery Occlusions in Noncontrast Brain Computed Tomography. &  &  &  & E \\
36 & Temporal-Spatial Feature Extraction of DSA Video and Its Application in AVM Diagnosis. &  &  &  & I \\
37 & Hunting for a coronary artery disease diagnosis in asymptomatic patients with diabetes mellitus: if, how and when. &  &  &  & E \\
38 & Automatic differentiation of ruptured and unruptured intracranial aneurysms on computed tomography angiography based on deep learning and radiomics. &  &  &  & E \\
39 & Improved U-Net for Guidewire Tip Segmentation in X-ray Fluoroscopy Images &  &  &  & I \\
40 & Prediction of Percutaneous Coronary Intervention Success in Patients With Moderate to Severe Coronary Artery Calcification Using Machine Learning Based on Coronary Angiography: Prospective Cohort Study. &  &  &  & I \\
41 & Multi-scale dual attention embedded U-shaped network for accurate segmentation of coronary vessels in digital subtraction angiography. &  &  &  & I \\
42 & Quantifying post-treatment vascular remodeling in brain aneurysms using WEKA-based machine learning: a pilot study. &  &  &  & I \\
43 & Use of biplane quantitative angiographic imaging with ensemble neural networks to assess reperfusion status during mechanical thrombectomy. &  &  &  & I \\
44 & Automatic Image Co-registration of Carotid Angiography and Intravascular Optical Coherence Tomography Based on Machine Learning Method: A Pilot Feasibility Study. &  &  &  & I \\
45 & A deep learning framework for segmentation and pose estimation of pedicle screw implants based on C-arm fluoroscopy. &  &  &  & I \\
46 & Attention residual convolution neural network based on U-net (AttentionResU-Net) for retina vessel segmentation &  &  &  & E \\
47 & Anomaly Detection in Optical Coherence Tomography Angiography (OCTA) with a Vector-Quantized Variational Auto-Encoder (VQ-VAE). &  &  &  & E \\
48 & Digital Subtraction Angiography Image Features under the Deep Learning Algorithm in Cardiovascular Interventional Treatment and Nursing for Vascular Restenosis. &  &  &  & I \\
49 & Cryo-balloon catheter localization in X-Ray fluoroscopy using U-net. &  &  &  & I \\
50 & Weakly-Supervised Convolutional Neural Networks for Vessel Segmentation in Cerebral Angiography &  &  &  & I \\
51 & Differences in retinal microvasculature between large artery atherosclerosis and small artery disease: an optical coherence tomography angiography study. &  &  &  & E \\
52 & A framework for intracranial aneurysm detection and rupture analysis on DSA. &  &  &  & I \\
53 & Fully automated unruptured intracranial aneurysm detection and segmentation from digital subtraction angiography series using an end-to-end spatiotemporal deep neural network &  &  &  & I \\
54 & Prediction of revascularization by coronary CT angiography using a machine learning ischemia risk score. &  &  &  & E \\
55 & Real-time Markerless Tracking of Lung Tumors based on 2-D Fluoroscopy Imaging using Convolutional LSTM. &  &  &  & I \\
56 & Enhancing percutaneous coronary intervention with heuristic path planning and deep-learning-based vascular segmentation. &  &  &  & I \\
57 & Automatic detection of pulmonary embolism in computed tomography pulmonary angiography using Scaled-YOLOv4. &  &  &  & E \\
58 & Noise reduction of three-dimensional rotational angiography (3DRA) images using residual encoder-decoder convolutional neural network (RED-CNN) &  &  &  & I \\
59 & Noninvasive Coronary Artery Disease Detection Using Retinal Images: A Multimodal Study. &  &  &  & E \\
60 & Radiomics-based assessment of ultra-widefield leakage patterns and vessel network architecture in the PERMEATE study: insights into treatment durability. &  &  &  & I \\
61 & Performance of a Deep-Learning Neural Network to Detect Intracranial Aneurysms from 3D TOF-MRA Compared to Human Readers. &  &  &  & E \\
62 & Auxiliary learning for anatomical localization of hierarchical vascular structures in digital subtraction angiography &  &  &  & I \\
63 & INTERMEDIATE AND DEEP CAPILLARY PLEXUSES IN MACHINE LEARNING SEGMENTATION OF HIGH-RESOLUTION OPTICAL COHERENCE TOMOGRAPHY IMAGING. &  &  &  & E \\
64 & Deep Learning-Based Digital Subtraction Angiography Characteristics in Nursing of Maintenance Hemodialysis Patients. &  &  &  & I \\
65 & Comparison of deep learning reconstruction algorithms to improve image quality of dual-energy carotid CT angiography under dual-low scan. &  &  &  & E \\
66 & The ApoB/ApoA-I ratio supersedes conventional lipids in predicting coronary artery disease and clinical phenotypes requiring revascularization. &  &  &  & E \\
67 & AutoFOX: An automated cross-modal 3D fusion framework of coronary X-ray angiography and OCT. &  &  &  & I \\
68 & Detection of Endoleak after Endovascular Aortic Repair through Deep Learning Based on Non-contrast CT. &  &  &  & E \\
69 & Normal-resolution vs. super-resolution deep learning reconstruction for diagnosis of functionally significant coronary stenosis using cardiac CT. &  &  &  & E \\
70 & Super-resolution deep learning reconstruction for improved quality of myocardial CT late enhancement. &  &  &  & E \\
71 & Image Quality Assessment of Deep Learning-Based Virtual Monoenergetic Images From Single-Energy CT Pulmonary Angiography. &  &  &  & E \\
72 & Coronary Artery Stent Evaluation by CTA: Impact of Deep Learning Reconstruction and Subtraction Technique. &  &  &  & E \\
73 & Fully Automatic Dual-Guidewire Segmentation for Coronary Bifurcation Lesion &  &  &  & I \\
74 & Questionnaire-based exposome-wide association studies (ExWAS) reveal expected and novel risk factors associated with cardiovascular outcomes in the Personalized Environment and Genes Study. &  &  &  & E \\
75 & Multitask Deep Learning for Automated Detection of Endoleak at Digital Subtraction Angiography during Endovascular Aneurysm Repair. &  &  &  & I \\
76 & What role for cardiac imaging in chronic coronary syndromes: review of the literature in light of the latest recommendations. &  &  &  & E \\
77 & Fully Automated Approaches for Localization of Intraoperative Electrocorticographic Electrodes &  &  &  & I \\
78 & Patch Neointima Technique in Acute Type A Aortic Dissection: Midterm Results of 147 Cases. &  &  &  & E \\
79 & Gastrointestinal bleeding detection on digital subtraction angiography using convolutional neural networks with and without temporal information. &  &  &  & I \\
80 & Automatic extraction and stenosis evaluation of coronary arteries in invasive coronary angiograms. &  &  &  & I \\
81 & Personalised pose estimation from single-plane moving fluoroscope images using deep convolutional neural networks. &  &  &  & I \\
82 & Deep Learning-based Angiogram Generation Model for Cerebral Angiography without Misregistration Artifacts. &  &  &  & I \\
83 & Direct Vascular Territory Segmentation on Cerebral Digital Subtraction Angiography &  &  &  & I \\
84 & Influence of diabetes mellitus on the diagnostic performance of machine learning-based coronary CT angiography-derived fractional flow reserve: a multicenter study. &  &  &  & E \\
85 & Enhancing Pulmonary Embolism Detection in COVID-19 Patients Through Advanced Deep Learning Techniques. &  &  &  & E \\
86 & Label-Efficient Data Augmentation with Video Diffusion Models for Guidewire Segmentation in Cardiac Fluoroscopy &  &  &  & I \\
87 & Establishment of machine learning-based tool for early detection of pulmonary embolism. &  &  &  & E \\
88 & Automated three-dimensional vessel reconstruction based on deep segmentation and bi-plane angiographic projections. &  &  &  & I \\
89 & A Dynamic Machine Learning Model to Predict Angiographic Vasospasm After Aneurysmal Subarachnoid Hemorrhage. &  &  &  & E \\
90 & Perfusion Maps Acquired From Dynamic Angiography MRI Using Deep Learning Approaches. &  &  &  & E \\
91 & Stroke prediction in elderly patients with atrial fibrillation using machine learning combined clinical and left atrial appendage imaging phenotypic features. &  &  &  & E \\
92 & Calculation of virtual 3D subtraction angiographies using conditional generative adversarial networks (cGANs). &  &  &  & I \\
93 & Validation Study Comparing Artificial Intelligence for Fully Automatic Aortic Aneurysms Segmentation and Diameter Measurements on Contrast and Noncontrast Enhanced Computed Tomography: A la mémoire de Lucie Salomon du Mont (In memory of Lucie Salomon du Mont). &  &  &  & E \\
94 & Phenotypic clustering of myocardial infarction complicated by out-of-hospital cardiac arrest using unsupervised machine learning. &  &  &  & E \\
95 & Automatic Pharyngeal Phase Recognition in Untrimmed Videofluoroscopic Swallowing Study Using Transfer Learning with Deep Convolutional Neural Networks. &  &  &  & I \\
96 & Efficient Quality Control with Mixed CT and CTA Datasets. &  &  &  & E \\
97 & Towards clinical translation of deep-learning based classification of DSA image sequences for stroke treatment &  &  &  & I \\
98 & AI-driven tongue image analysis for diagnosing and predicting coronary artery disease. &  &  &  & E \\
99 & Comparison of Image Quality Reconstructed Using Iterative Reconstruction and Deep Learning Algorithms Under Varying Dose Reductions in Dual-Energy Carotid CT Angiography. &  &  &  & E \\
100 & GaraMoSt: Parallel Multi-Granularity Motion and Structural Modeling for Efficient Multi-Frame Interpolation in DSA Images &  &  &  & I \\
101 & Fluoroscopic Image Denoising with Feature Preserving Residual Noise Learning &  &  &  & I \\
102 & TiAVox: Time-aware Attenuation Voxels for Sparse-view 4D DSA Reconstruction &  &  &  & I \\
103 & Evaluation of Optic Disc, Retinal Vascular Structures, and Acircularity Index in Patients with Idiopathic Macular Telangiectasia Type 2. &  &  &  & E \\
104 & Development and validation of a neural network survival prediction model for ischemic heart disease. &  &  &  & E \\
105 & Studies of Plausibility Using Deep Learning Techniques to Automatically Detect Retinal Neovascularization in Enface Optical Coherence Tomography Angiography Images &  &  &  & E \\
106 & Enhancing precision in aneurysm volume measurement: A comparative study of techniques including an artificial intelligence-based method for endovascular coiling. &  &  &  & E \\
107 & Uncertain Feature-refinement Attention Unet: Considering Suitable Convolutional Neural Network Model for Real-time Segmentation in Markerless Tumor Tracking. &  &  &  & I \\
108 & Radiogenomics and Artificial Intelligence Approaches Applied to Cardiac Computed Tomography Angiography and Cardiac Magnetic Resonance for Precision Medicine in Coronary Heart Disease: A Systematic Review. &  &  &  & E \\
109 & Energy-based Kinetic Energy Features for Online Dynamic Security Assessment using Bayes by Backprop Machine Learning Algorithm &  &  &  & E \\
110 & Performance of deep learning-based autodetection of arterial stenosis on head and neck CT angiography: an independent external validation study. &  &  &  & E \\
111 & Classifying Ruptured Middle Cerebral Artery Aneurysms With a Machine Learning Based, Radiomics-Morphological Model: A Multicentral Study. &  &  &  & E \\
112 & Autonomous Guidewire Navigation in Dynamic Environments &  &  &  & I \\
113 & Deep learning for estimating right ventricular function from routine coronary angiography. &  &  &  & I \\
114 & Automatic image-based tracking of gadolinium-filled balloon wedge catheters for MRI-guided cardiac catheterization using deep learning. &  &  &  & E \\
115 & Background Subtraction Angiography with Deep Learning Using Multi-frame Spatiotemporal Angiographic Input. &  &  &  & I \\
116 & Development of a machine learning-based real-time location system to streamline acute endovascular intervention in acute stroke: a proof-of-concept study. &  &  &  & E \\
117 & Multi-dimensional cascades neural network models for the segmentation of retinal vessels in colour fundus images &  &  &  & E \\
118 & AI Quantification of Vascular Lesions in Mouse Fundus Fluorescein Angiography. &  &  &  & I \\
119 & The Clinical Utility Of AI Models In Diagnostic And Endovascular Radiology &  &  &  & E \\
120 & Deep reinforcement learning for cerebral anterior vessel tree extraction from 3D CTA images. &  &  &  & E \\
121 & Benchmark dataset for clot detection in ischemic stroke vessel-based imaging: CODEC-IV. &  &  &  & E \\
122 & Proximal femur segmentation and quantification in dual-energy subtraction tomosynthesis: A novel approach to fracture risk assessment. &  &  &  & I \\
123 & Femoropopliteal Artery Stent Detection in Fluoroscopy Using Otsu's Preprocessing and CNN Segmentation &  &  &  & I \\
124 & Automatic viewing angles prediction for endovascular treatment in intracranial aneurysm &  &  &  & I \\
125 & Deep Learning-Based Automated Thrombolysis in Cerebral Infarction Scoring: A Timely Proof-of-Principle Study. &  &  &  & I \\
126 & Stent-RD-UNet: A Deep Learning Network for Real-Time Pose Estimation and Tracking of Aortic Stents in Image-Assisted Interventions &  &  &  & I \\
127 & Automated catheter segmentation and tip detection in cerebral angiography with topology-aware geometric deep learning. &  &  &  & I \\
128 & Optimizing Input Selection for Cardiac Model Training and Inference: An Efficient 3D CNN-based Approach to Automate Coronary Angiogram Video Selection &  &  &  & I \\
129 & Prediction of treatment outcome for branch retinal vein occlusion using convolutional neural network-based retinal fluorescein angiography. &  &  &  & I \\
130 & Deep learning based tracked X-ray for surgery guidance &  &  &  & I \\
131 & Intraoperative guidance of orthopaedic instruments using 3D correspondence of 2D object instance segmentations &  &  &  & I \\
132 & DSA Image Analysis of Clinical Features and Nursing Care of Cerebral Aneurysm Patients Based on the Deep Learning Algorithm. &  &  &  & I \\
133 & Deep Learning-Based Approach to Automatically Assess Coronary Distensibility Following Kawasaki Disease. &  &  &  & I \\
134 & Editor's choice to the May 2022 issue : Fully automated quantification of cardiac chambers and function in 2D echo by Deep Learning, and a modern atlas of invasive coronary angiographic views. &  &  &  & E \\
135 & Lesion detection in age-related macular degeneration with a multi-modal imaging and machine learning approach. &  &  &  & E \\
136 & AI-quantified epicardial adipose tissue and prediction of future myocardial infarction in patients with cardiometabolic disease: a post-hoc analysis from the SCOT-HEART trial. &  &  &  & E \\
137 & Dynamic-Computed Tomography Angiography for Cerebral Vessel Templates and Segmentation &  &  &  & E \\
138 & Vision-language model for report generation and outcome prediction in CT pulmonary angiogram. &  &  &  & E \\
139 & Automated Echocardiographic Detection of Severe Coronary Artery Disease Using Artificial Intelligence. &  &  &  & E \\
140 & A Memory Efficient Adversarial Attention Tree-Structured Deep Learning Model for Classification &  &  &  & E \\
141 & Deep Learning Analysis of Indocyanine Green Fluoroscopy of Ureters in Robotic Cystectomy: Toward Reducing Ureteroenteric Strictures. &  &  &  & I \\
142 & SPARNet: A Framework for Airway Invasion Tracking from Fluoroscopic Videos of Dysphagia Patients. &  &  &  & I \\
143 & Automated assessment of right ventricular systolic function from coronary angiograms with video-based artificial intelligence algorithms: development, validation, comparison against humans, and prospective deployment. &  &  &  & I \\
144 & S-MEOD: A Novel Evaluation Metric for Frame-Based Medical Object Detection. &  &  &  & I \\
145 & Fluoroscopic image-driven deep learning model for predicting intussusception irreducibility during air enema in children. &  &  &  & I \\
146 & YOLO11-based detection of manometry sensors in video-fluoroscopy imaging for computer-aided multimodal assessment of swallowing. &  &  &  & I \\
147 & Toward real-time autonomous navigation: transformer-based catheter tip tracking in fluoroscopy. &  &  &  & I \\
148 & AngioCAD: A public x-ray angiography dataset and an adaptive fusion framework for stenosis detection. &  &  &  & I \\
149 & TransUNet for automated segmentation of intracranial aneurysms in digital subtraction angiography: A dual-center validation study. &  &  &  & I \\
150 & Quantitative Assessment of Fluorescein Angiography Leakage via Deep Learning in Pediatric Uveitis: Correlation with Clinical Parameters. &  &  &  & I \\
151 & A Deep Learning Model for Artery of Adamkiewicz Detection in Bronchial Artery Embolization: A Multicenter Retrospective Study. &  &  &  & I \\
152 & PASAformer: Cerebrovascular Disease Classification with Medical Prior-Guided Adapter and Pathology-Aware Sparse Attention. &  &  &  & I \\
153 & Artificial intelligence risk stratification from dynamic digital subtraction angiography radiomics predicts pulmonary embolism and associates with clinical outcomes in deep vein thrombosis: A retrospective cohort study. &  &  &  & I \\
154 & Machine learning-enabled prediction of hemorrhagic transformation post-thrombectomy using quantitative DSA. &  &  &  & I \\
155 & Artificial intelligence classification of rectal neoplasia by endoscopic fluorescence perfusion analysis. &  &  &  & I \\
156 & Identifying Pathogenesis of Acute Coronary Syndromes using Sequence Contrastive Learning in Coronary Angiography. &  &  &  & I \\
157 & Multi-graph Graph matching for coronary artery semantic labeling in invasive coronary angiograms. &  &  &  & I \\
158 & Catheter detection and segmentation in X-ray images via multi-task learning. &  &  &  & I \\
159 & Deep learning-based frame synthesis enables radiation dose reduction in digital subtraction angiography imaging: a multicenter study &  &  &  & I \\
160 & CDSA-Net:Collaborative Decoupling of Vascular Structure and Background for High-Fidelity Coronary Digital Subtraction Angiography &  &  &  & I \\
161 & Real-Time Tracking of Flexible Markers in Low-Contrast Fluoroscopy Using a Deep Neural Network Trained Solely on Synthetic Data &  &  &  & I \\
162 & SIX: A Self-Supervised Interactive Algorithm for Low-Dose X-Ray Fluoroscopy Image Denoising &  &  &  & I \\
163 & Coronary Artery Segmentation and Vessel-Type Classification in X-Ray Angiography &  &  &  & I \\
164 & From Prompts to Pipelines: Evaluating LLM-Generated Medical Image Segmentation Baselines &  &  &  & I \\
165 & Development and Validation of an Interpretable Hemodynamics-Based Machine Learning Model for Predicting Cerebral Arteriovenous Malformation Rupture. &  &  &  & I \\
166 & Texture-Based Preprocessing Framework with nnU-Net Model for Accurate Intracranial Artery Segmentation. &  &  &  & I \\
167 & Multimodal machine learning for 5-year mortality prediction after percutaneous coronary intervention. &  &  &  & I \\
168 & AI-Based Denoising and Identification of Neuromodulation Lead Locations from Limited Intraoperative Fluoroscopy Images. &  &  &  & I \\
169 & Optimized aortic root segmentation during transcatheter aortic valve implantation. &  &  &  & I \\
170 & Marker-Less Lung Tumor Tracking from Real-Time Color X-Ray Fluoroscopic Images Using Cross-Patient Deep Learning Model. &  &  &  & I \\
171 & Feasibility of Multimodal Deep Learning for Automated Staging of Familial Exudative Vitreoretinopathy Using Color Fundus Photographs and Fluorescein Angiography. &  &  &  & I \\
172 & Intracranial aneurysm segmentation on digital subtraction angiography: a retrospective and multi-center study. &  &  &  & I \\
173 & Real-Time Artificial Intelligence-Assisted Middle Meningeal Artery Embolization Using Liquid Embolic Agents for Chronic Subdural Hematoma: A Preliminary Experience. &  &  &  & I \\
174 & Quantitative DSA Analysis of MCA Aneurysms Using SymDIRECT Pixel Clustering: A Novel Framework for Objective Post-Treatment Evaluation. &  &  &  & I \\
175 & Deep learning-based dual-energy subtraction synthesis from single-energy kV x-ray fluoroscopy for markerless tumor tracking. &  &  &  & I \\
176 & Explainable AI for raising confidence in deep learning-based tumor tracking models. &  &  &  & I \\
177 & Design and optimization of an automatic deep learning-based cerebral reperfusion scoring (TICI) using thrombus localization. &  &  &  & I \\
178 & Semi-automatic segmentation of elongated interventional instruments for online calibration of C-arm imaging system. &  &  &  & I \\
179 & Generative deep-learning-model based contrast enhancement for digital subtraction angiography using a text-conditioned image-to-image model. &  &  &  & I \\
180 & Machine Learning-Assisted Prediction of Persistent Incomplete Occlusion in Intracranial Aneurysms From Angiographic Parametric Imaging-Derived Features. &  &  &  & I \\
181 & End-to-end 2D/3D registration from pre-operative MRI to intra-operative fluoroscopy for orthopedic procedures. &  &  &  & I \\
182 & Real-Time Coronary Artery Dominance Classification from Angiographic Images Using Advanced Deep Video Architectures. &  &  &  & I \\
183 & Training a deep learning model to predict the anatomy irradiated in fluoroscopic x-ray images. &  &  &  & I \\
184 & Improving automatic cerebral 3D-2D CTA-DSA registration. &  &  &  & I \\
185 & Automatic transformer-based grading of multiple retinal inflammatory signs in uveitis on fluorescein angiography. &  &  &  & I \\
186 & Accuracy of segment anything model for classification of vascular stenosis in digital subtraction angiography. &  &  &  & I \\
187 & Lightweight attention network for guidewire segmentation and localization in clinical fluoroscopic images of vascular interventional surgery. &  &  &  & I \\
188 & perfDSA: Automatic Perfusion Imaging in Cerebral Digital Subtraction Angiography. &  &  &  & I \\
189 & Optimizing Input Selection for Cardiac Model Training and Inference: An Efficient 3D Convolutional Neural Networks-Based Approach to Automate Coronary Angiogram Video Selection. &  &  &  & I \\
190 & Linking sequence restoration capability of shuffled coronary angiography to coronary artery disease diagnosis. &  &  &  & I \\
191 & Machine Learning Analysis of Videourodynamics to Predict Incident Hydronephrosis in Patients With Spina Bifida. &  &  &  & I \\
192 & Intraoperative stenosis detection in X-ray coronary angiography via temporal fusion and attention-based CNN. &  &  &  & I \\
193 & CVFSNet: A Cross View Fusion Scoring Network for end-to-end mTICI scoring. &  &  &  & I \\
194 & CDKD-w+: A Keyframe Recognition Method for Coronary Digital Subtraction Angiography Video Sequence Based on w+ Space Encoding. &  &  &  & I \\
195 & Technique selection and technical developments for 2D dual-energy subtraction angiography on an interventional C-arm. &  &  &  & I \\
196 & Development of a Deep Learning Model for Automating Implant Position in Total Hip Arthroplasty. &  &  &  & I \\
197 & Automated stenosis estimation of coronary angiographies using end-to-end learning. &  &  &  & I \\
198 & Deep Learning Models for Automatic Classification of Anatomic Location in Abdominopelvic Digital Subtraction Angiography. &  &  &  & I \\
199 & Deep Learning-Based Detection of Malignant Bile Duct Stenosis in Fluoroscopy Images of Endoscopic Retrograde Cholangiopancreatography. &  &  &  & I \\
200 & Rupture risk assessment in cerebral arteriovenous malformations: an ensemble model using hemodynamic and morphological features. &  &  &  & I \\
201 & Large-scale pretrained frame generative model enables real-time low-dose DSA imaging: An AI system development and multi-center validation study. &  &  &  & I \\
202 & Assessing superficial temporal artery-middle cerebral artery anastomosis patency using FLOW 800 hemodynamics. &  &  &  & I \\
203 & Predicting Intracranial Aneurysm Rupture: A Multifactor Analysis Combining Radscore, Morphology, and PHASES Parameters. &  &  &  & I \\
204 & Enhanced Cerebrovascular Extraction Using Vessel-Specific Preprocessing of Time-Series Digital Subtraction Angiograph &  &  &  & I \\
205 & OccluNet: Spatio-Temporal Deep Learning for Occlusion Detection on DSA &  &  &  & I \\
206 & CTNet: A CNN-Transformer Dual-Branch Network for 2D-3D Rigid Registration &  &  &  & I \\
207 & Deep learning fluoroscopy denoising: a beam quality insensitive approach &  &  &  & I \\
208 & DSA-NRP: No-Reflow Prediction from Angiographic Perfusion Dynamics in Stroke EVT &  &  &  & I \\
209 & Abstract 097: Real‐time AI‐assisted MMA embolization using liquid embolic agents for chronic subdural hematoma &  &  &  & I \\
210 & Reducing motion artifacts in craniocervical background subtraction angiography with deformable registration and unsupervised deep learning. &  &  &  & I \\
211 & Optimizing Modified Barium Swallow Exam Workflow: Automating Pre-Analysis Video Sorting in Swallowing Function Assessment. &  &  &  & I \\
212 & Parathyroid gland identification and angiography classification using simple machine learning methods. &  &  &  & I \\
213 & First In-Vivo demonstration of 1000fps High Speed Coronary Angiography (HSCA) in a swine animal model. &  &  &  & I \\
214 & Simulation-informed learning for time-resolved angiographic contrast agent concentration reconstruction. &  &  &  & I \\
215 & DSA Quantitative Analysis and Predictive Modeling of Obliteration in Cerebral AVM following Stereotactic Radiosurgery. &  &  &  & I \\
216 & Automatic localization of anatomical landmarks in head cine fluoroscopy images via deep learning. &  &  &  & I \\
217 & Optimizing neurointerventional procedures: an algorithm for embolization coil detection and automated collimation to enable dose reduction. &  &  &  & I \\
218 & Automated Laryngeal Invasion Detector of Boluses in Videofluoroscopic Swallowing Study Videos Using Action Recognition-Based Networks. &  &  &  & I \\
219 & A novel post-processing technique for correcting symmetric implant ambiguity in measuring total knee arthroplasty kinematics from single-plane fluoroscopy. &  &  &  & I \\
220 & CAVE: Cerebral artery-vein segmentation in digital subtraction angiography. &  &  &  & I \\
221 & A deep learning method to identify and localize large-vessel occlusions from cerebral digital subtraction angiography. &  &  &  & I \\
222 & A feasibility study assessing quantitative indocyanine green angiographic predictors of reconstructive complications following nipple-sparing mastectomy. &  &  &  & I \\
223 & Cross-modality Labeling Enables Noninvasive Capillary Quantification as a Sensitive Biomarker for Assessing Cardiovascular Risk. &  &  &  & I \\
224 & Automatic Detection of 30 Fundus Diseases Using Ultra-Widefield Fluorescein Angiography with Deep Experts Aggregation. &  &  &  & I \\
225 & Significantly better precision with new machine learning versus manual image registration software in processing images from single-plane fluoroscopy to determine tibiofemoral kinematics following total knee replacement. &  &  &  & I \\
226 & Radiation reduction for interventional radiology imaging: a video frame interpolation solution. &  &  &  & I \\
227 & Training of a deep learning based digital subtraction angiography method using synthetic data. &  &  &  & I \\
228 & Topology observing 3D device reconstruction from continuous-sweep limited angle fluoroscopy. &  &  &  & I \\
229 & Development of an Artificial Intelligence Tool for Intraoperative Guidance During Endovascular Abdominal Aortic Aneurysm Repair. &  &  &  & I \\
230 & Fully automatic algorithm for detecting and tracking anatomical shoulder landmarks on fluoroscopy images with artificial intelligence. &  &  &  & I \\
231 & A Fully Differentiable Framework for 2D/3D Registration and the Projective Spatial Transformers. &  &  &  & I \\
232 & Deep learning-based cerebral aneurysm segmentation and morphological analysis with three-dimensional rotational angiography. &  &  &  & I \\
233 & Unsupervised deep learning for the matching of vascular anatomy in multiple digital subtraction angiograms &  &  &  & I \\
234 & DFA-Net: Dual multi-scale feature aggregation network for vessel segmentation in X-ray digital subtraction angiography &  &  &  & I \\
235 & Artifact-aware Digital Subtraction Angiogram Image Generation for Head and Neck Vessels &  &  &  & I \\
236 & Enhancing texture detail recovery in low-dose x-ray fluoroscopic images with a multi-frame deep learning framework &  &  &  & I \\
237 & ABLSpineLevelCheck: Localization of Vertebral Levels on Fluoroscopy via Semi-supervised Abductive Learning &  &  &  & I \\
238 & A Patient-Specific Registration of Coronary Angiogram-Fluoroscopy by Similarity- Based Transfer Learning &  &  &  & I \\
239 & Use of Yolo Detection for 3D Pose Tracking of Cardiac Catheters Using Bi-Plane Fluoroscopy &  &  &  & I \\
240 & A study on the influence of radiation dose on device tracking accuracy during continuous-sweep limited angle fluoroscopy &  &  &  & I \\
241 & Unsupervised Training of a Dynamic Context-Aware Deep Denoising Framework for Low-Dose Fluoroscopic Imaging &  &  &  & I \\
242 & Deformable Model-to-Image Registration Toward Augmented Reality-Guided Endovascular Interventions &  &  &  & I \\
243 & 4DRGS: 4D Radiative Gaussian Splatting for Efficient 3D Vessel Reconstruction from Sparse-View Dynamic DSA Images &  &  &  & I \\
244 & MoSt-DSA: Modeling Motion and Structural Interactions for Direct Multi-Frame Interpolation in DSA Images &  &  &  & I \\
245 & A Deep Learning Model for Coronary Artery Segmentation and Quantitative Stenosis Detection in Angiographic Images &  &  &  & I \\
246 & Deep learning approaches for image restoration in invasive coronary angiography &  &  &  & I \\
247 & Automated Interpretation of Fundus Fluorescein Angiography with Multi-Retinal Lesion Segmentation &  &  &  & I \\
248 & Deep Learning to Estimate Left Ventricular Ejection Fraction From Routine Coronary Angiographic Images. &  &  &  & I \\
249 & Pelphix: Surgical Phase Recognition from X-ray Images in Percutaneous Pelvic Fixation. &  &  &  & I \\
250 & USING CONVOLUTIONAL NEURAL NETWORK-BASED SEGMENTATION FOR IMAGE-BASED COMPUTATIONAL FLUID DYNAMICS SIMULATIONS OF BRAIN ANEURYSMS: INITIAL EXPERIENCE IN AUTOMATED MODEL CREATION. &  &  &  & I \\
251 & An Improved Microaneurysm Detection Model Based on SwinIR and YOLOv8. &  &  &  & I \\
252 & TBENet:A two-branch boundary enhancement Network for cerebrovascular segmentation. &  &  &  & I \\
253 & Use of deep learning to segment bolus during videofluoroscopic swallow studies. &  &  &  & I \\
254 & A framework of myocardial bridge detection with x-ray angiography sequence. &  &  &  & I \\
255 & Joint Track Machine Learning: An Autonomous Method of Measuring Total Knee Arthroplasty Kinematics From Single-Plane X-Ray Images. &  &  &  & I \\
256 & Automatic prediction model for online diaphragm motion tracking based on optical surface monitoring by machine learning. &  &  &  & I \\
257 & Deep Learning of Videourodynamics to Classify Bladder Dysfunction Severity in Patients With Spina Bifida. &  &  &  & I \\
258 & High-throughput 3DRA segmentation of brain vasculature and aneurysms using deep learning. &  &  &  & I \\
259 & A deep learning based dual encoder-decoder framework for anatomical structure segmentation in chest X-ray images. &  &  &  & I \\
260 & Coronary artery stenosis detection via proposal-shifted spatial-temporal transformer in X-ray angiography. &  &  &  & I \\
261 & Deep Learning Subtraction Angiography: Improved Generalizability with Transfer Learning. &  &  &  & I \\
262 & Jigsaw training-based background reverse attention transformer network for guidewire segmentation. &  &  &  & I \\
263 & Automated detection of arterial landmarks and vascular occlusions in patients with acute stroke receiving digital subtraction angiography using deep learning. &  &  &  & I \\
264 & Learning spatiotemporal features of DSA using 3D CNN and BiConvGRU for ischemic moyamoya disease detection. &  &  &  & I \\
265 & A Coordinate-Regression-Based Deep Learning Model for Catheter Detection during Structural Heart Interventions &  &  &  & I \\
266 & 3D Guidewire Shape Reconstruction from Monoplane Fluoroscopic Images &  &  &  & I \\
267 & ConTrack: Contextual Transformer for Device Tracking in X-ray &  &  &  & I \\
268 & A motion-level-aware denoising framework for x-ray fluoroscopic images &  &  &  & I \\
269 & Prolonged turbulence and deposition of LDL cholesterol caused atherosclerotic lesion in femoral arteries and spared the iliac arteries. Why? A dynamic angiography and machine learning analysis &  &  &  & I \\
270 & Machine learning-assisted fluoroscopy of bladder function in awake mice. &  &  &  & I \\
271 & Automatic landmark detection and mapping for 2D/3D registration with BoneNet. &  &  &  & I \\
272 & Initial investigation of the use of angiographic parametric imaging for early prognosis of delayed cerebral ischemia in patients with subarachnoid hemorrhage. &  &  &  & I \\
273 & Quantitative angiography prognosis of intracranial aneurysm treatment failure using parametric imaging and distal vessel analysis. &  &  &  & I \\
274 & Real-time Detection of Patient Head Position and Cephalometric Landmarks from Neuro-Interventional Procedure Images Using Machine Learning for Patient Eye-Lens Dose Prediction. &  &  &  & I \\
275 & Automated Detection of Vascular Leakage in Fluorescein Angiography - A Proof of Concept. &  &  &  & I \\
276 & InterNet: Detection of Active Abdominal Arterial Bleeding Using Emergency Digital Subtraction Angiography Imaging With Two-Stage Deep Learning. &  &  &  & I \\
277 & Optical flow estimation of coronary angiography sequences based on semi-supervised learning. &  &  &  & I \\
278 & Deep learning-based classification of DSA image sequences of patients with acute ischemic stroke. &  &  &  & I \\
279 & Deep learning-based framework for motion-compensated image fusion in catheterization procedures. &  &  &  & I \\
280 & Real-time automatic prediction of treatment response to transcatheter arterial chemoembolization in patients with hepatocellular carcinoma using deep learning based on digital subtraction angiography videos. &  &  &  & I \\
281 & Deep Learning Segmentation in 2D X-ray Images and Non-Rigid Registration in Multi-Modality Images of Coronary Arteries. &  &  &  & I \\
282 & Maskless 2-Dimensional Digital Subtraction Angiography Generation Model for Abdominal Vasculature using Deep Learning. &  &  &  & I \\
283 & Using a Convolutional Neural Network and Convolutional Long Short-term Memory to Automatically Detect Aneurysms on 2D Digital Subtraction Angiography Images: Framework Development and Validation. &  &  &  & I \\
284 & Spatio-temporal deep learning for automatic detection of intracranial vessel perforation in digital subtraction angiography during endovascular thrombectomy. &  &  &  & I \\
285 & Edge-enhancement densenet for X-ray fluoroscopy image denoising in cardiac electrophysiology procedures. &  &  &  & I \\
286 & Identification of successful cerebral reperfusions (mTICI ≥2b) using an artificial intelligence strategy. &  &  &  & I \\
287 & A Unified Framework for Multi-Guidewire Endpoint Localization in Fluoroscopy Images. &  &  &  & I \\
288 & Deep learning-enabled ultra-widefield retinal vessel segmentation with an automated quality-optimized angiographic phase selection tool. &  &  &  & I \\
289 & Personalized Pre- and Post-Operative Hemodynamic Assessment of Aortic Coarctation from 3D Rotational Angiography. &  &  &  & I \\
290 & Robust Landmark-based Stent Tracking in X-ray Fluoroscopy &  &  &  & I \\
291 & [Development of the DSA Method for Coronary Angiography Using Deep Learning Techniques]. &  &  &  & I \\
292 & Breathing-Compensated Neural Networks for Real Time C-Arm Pose Estimation in Lung CT-Fluoroscopy Registration &  &  &  & I \\
293 & Towards Real-time pose estimation of the Mitral Valve Robot under C-arm X-ray Fluoroscopy. &  &  &  & I \\
294 & Towards real time guide wire shape extraction in fluoroscopic sequences: A two phase deep learning scheme to extract sparse curvilinear structures. &  &  &  & I \\
295 & Multi-Modality guidance based surgical navigation for percutaneous endoscopic transforaminal discectomy. &  &  &  & I \\
296 & Deformable 3D-2D registration for high-precision guidance and verification of neuroelectrode placement. &  &  &  & I \\
297 & Transfer learning from an artificial radiograph-landmark dataset for registration of the anatomic skull model to dual fluoroscopic X-ray images. &  &  &  & I \\
298 & Automated Grading of Diabetic Retinopathy with Ultra-Widefield Fluorescein Angiography and Deep Learning. &  &  &  & I \\
299 & Deep learning-based reconstruction of interventional tools and devices from four X-ray projections for tomographic interventional guidance. &  &  &  & I \\
300 & Robust Orthogonal-View 2-D/3-D Rigid Registration for Minimally Invasive Surgery. &  &  &  & I \\
301 & The Aneurysm Occlusion Assistant, an AI platform for real time surgical guidance of intracranial aneurysms. &  &  &  & I \\
302 & autoTICI: Automatic Brain Tissue Reperfusion Scoring on 2D DSA Images of Acute Ischemic Stroke Patients. &  &  &  & I \\
303 & Automatic detection of leakage point in central serous chorioretinopathy of fundus fluorescein angiography based on time sequence deep learning. &  &  &  & I \\
304 & Multi-View Convolutional Neural Networks in Rupture Risk Assessment of Small, Unruptured Intracranial Aneurysms. &  &  &  & I \\
305 & Deep neural network-based detection and segmentation of intracranial aneurysms on 3D rotational DSA. &  &  &  & I \\
306 & Use of quantitative angiographic methods with a data-driven model to evaluate reperfusion status (mTICI) during thrombectomy. &  &  &  & I \\
307 & Transfer Learning for Nonrigid 2D/3D Cardiovascular Images Registration. &  &  &  & I \\
308 & Recognition of moyamoya disease and its hemorrhagic risk using deep learning algorithms: sourced from retrospective studies. &  &  &  & I \\
309 & Quantitative analysis of flow vortices: differentiation of unruptured and ruptured medium-sized middle cerebral artery aneurysms. &  &  &  & I \\
310 & Abstract P494: Automatically Predicting Modified Treatment in Cerebral Ischemia Scores From Patient Digital Subtraction Angiography Using Deep Learning &  &  &  & I \\
311 & Model-to-Image Registration via Deep Learning towards Image-Guided Endovascular Interventions &  &  &  & I \\
312 & Deep-Learning-Based Registration of Diagnostic Angiogram and Live Fluoroscopy for Percutaneous Coronary Intervention &  &  &  & I \\
313 & Stent enhancement using marker response network and focus conversion learning in x-ray fluoroscopy sequences &  &  &  & I \\
314 & Using a Convolutional Neural Network and Convolutional Long Short-term Memory to Automatically Detect Aneurysms on 2D-DSA Images (Preprint) &  &  &  & I \\
315 & Automatic detection and tracking of the region of interest during fluoroscopy-guided procedures for radiation exposure reduction &  &  &  & I \\
316 & Real-Time Multi-Guidewire Endpoint Localization in Fluoroscopy Images &  &  &  & I \\
317 & Automated Deep Learning Analysis of Angiography Video Sequences for Coronary Artery Disease &  &  &  & I \\
318 & End-to-End Deep Learning Model for Cardiac Cycle Synchronization from Multi-View Angiographic Sequences. &  &  &  & I \\
319 & Deep neural networks for ECG-free cardiac phase and end-diastolic frame detection on coronary angiographies. &  &  &  & I \\
320 & Pyramid attention recurrent networks for real-time guidewire segmentation and tracking in intraoperative X-ray fluoroscopy. &  &  &  & I \\
321 & Catheter segmentation in X-ray fluoroscopy using synthetic data and transfer learning with light U-nets. &  &  &  & I \\
322 & A machine learning-based real-time tumor tracking system for fluoroscopic gating of lung radiotherapy. &  &  &  & I \\
323 & Adaptive weighted log subtraction based on neural networks for markerless tumor tracking using dual-energy fluoroscopy. &  &  &  & I \\
324 & Automatic Diagnosis Based on Spatial Information Fusion Feature for Intracranial Aneurysm. &  &  &  & I \\
325 & Automatic radiomic feature extraction using deep learning for angiographic parametric imaging of intracranial aneurysms. &  &  &  & I \\
326 & Semantic Segmentation to Extract Coronary Arteries in Fluoroscopy Angiograms &  &  &  & I \\
327 & Multiscale dense convolutional neural network for DSA cerebrovascular segmentation &  &  &  & I \\
328 & Fully automated intracranial aneurysm detection and segmentation from digital subtraction angiography series using an end-to-end spatiotemporal deep neural network &  &  &  & I \\
329 & Optimization of DSA image data input to a machine learning aneurysm identifier &  &  &  & I \\
330 & Image processing using Convolutional Neural Network (CNN) for Region of Interest (ROI) fluoroscopy &  &  &  & I \\
331 & Multi-resolution CNN for brain vessel segmentation from cerebrovascular images of intracranial aneurysm: a comparison of U-Net and DeepMedic &  &  &  & I \\
332 & Deep learning convolutional networks for image quality assessment in ultra-widefield fluorescein angiography &  &  &  & I \\
333 & Computer-assisted intra-operative verification of surgical outcome for the treatment of syndesmotic injuries through contralateral side comparison. &  &  &  & I \\
334 & Enabling machine learning in X-ray-based procedures via realistic simulation of image formation. &  &  &  & I \\
335 & A deep learning framework for automatic detection of arbitrarily shaped fiducial markers in intrafraction fluoroscopic images. &  &  &  & I \\
336 & Commissioning of a fluoroscopic-based real-time markerless tumor tracking system in a superconducting rotating gantry for carbon-ion pencil beam scanning treatment. &  &  &  & I \\
337 & Deep learning based guidewire segmentation in x-ray images &  &  &  & I \\
338 & Deformable Cardiovascular Image Registration via Multi-Channel Convolutional Neural Network &  &  &  & I \\
339 & Use of a convolutional neural network for aneurysm identification in digital subtraction angiography &  &  &  & I \\
340 & Feasibility study for use of angiographic parametric imaging and deep neural networks for intracranial aneurysm occlusion prediction &  &  &  & I \\
341 & Feasibility study of deep neural networks to classify intracranial aneurysms using angiographic parametric imaging &  &  &  & I \\
342 & Microcatheter Shaping for Intracranial Aneurysm Coiling Using Artificial Intelligence-Assisted Technique &  &  &  & I \\
343 & Automatic Annotation of Hip Anatomy in Fluoroscopy for Robust and Efficient 2D/3D Registration &  &  &  & I \\
344 & Estimation of attenuator mask from region of interest (ROI) dose-reduced images for brightness equalization using convolutional neural networks &  &  &  & I \\
345 & A novel real-time computational framework for detecting catheters and rigid guidewires in cardiac catheterization procedures. &  &  &  & I \\
346 & Automatic hyoid bone detection in fluoroscopic images using deep learning. &  &  &  & I \\
347 & Comparison of image quality of 40 keV virtual monoenergetic images of vertebral arteries using DLIR and ASIR-V algorithms under a dual-low scanning protocol. &  &  &  & E \\
348 & Guiding Preventive Care Strategies for Patients With Atherosclerotic Plaque on Coronary CTA: Primary Outcomes of the DECIDE Registry. &  &  &  & E \\
349 & Analysis of emerging trends in artificial intelligence for education in Nigeria &  &  &  & E \\
350 & Cardiac CT: Competition, complimentary or confounder. &  &  &  & E \\
351 & Quantitative Optical Coherence Tomography Detection of New Neovascular Branches and Association With Exudation Recurrence in Neovascular Age-Related Macular Degeneration. &  &  &  & E \\
352 & Diagnostic performance of deep learning-based coronary computed tomography angiography in detecting coronary artery stenosis. &  &  &  & E \\
353 & Artificial Intelligence in Cardiac Electrophysiology: Enhancing Mapping and Ablation Precision. &  &  &  & E \\
354 & The Features of Vessel Densities and Hemorheological Parameters in Patients With Central Retinal Vein Occlusion: A Prospective Randomized Case-Control Study. &  &  &  & E \\
355 & Update on Optical Coherence Tomography and Optical Coherence Tomography Angiography Imaging in Proliferative Diabetic Retinopathy. &  &  &  & E \\
356 & Artificial Intelligence Applications in Sickle Cell Retinopathy Imaging: Current Progress, Challenges, and Future Directions. &  &  &  & E \\
357 & Importance of coronary artery lumen size in the relationship between coronary artery plaque and vessel-specific ischemia: A post hoc analysis of CREDENCE and PACIFIC-1. &  &  &  & E \\
358 & A Cross-Sectional Study Assessing the Suitability of ChatGPT and DeepSeek AI for Generating Patient Education Guides on Imaging Modalities in Stroke. &  &  &  & E \\
359 & Multi-Plexus Nonperfusion Area Segmentation in Widefield OCT Angiography Using a Deep Convolutional Neural Network. &  &  &  & E \\
360 & Evaluating Deep Learning Techniques for Detecting Aneurysmal Subarachnoid Hemorrhage: A Comparative Analysis of Convolutional Neural Network and Transfer Learning Models. &  &  &  & E \\
361 & Association Between Flow Acceleration in the Carotid Artery and Intracranial Aneurysms. &  &  &  & E \\
362 & Accuracy of GPT-5 and GPT-4o in diagnosing STEMI from 12-Lead ECGs: A comparative study with cardiologists and emergency physicians. &  &  &  & E \\
363 & External Validation of Fully-Automated Infrarenal Maximum Aortic Aneurysm Diameter Measurements in Computed Tomography Angiography Scans Using Artificial Intelligence (PRAEVAorta 2). &  &  &  & E \\
364 & Quantitative Assessment of Three-Dimensional Choroidal Vascularity and Choriocapillaris Flow Signal Voids in Myopic Patients Using SS-OCTA. &  &  &  & E \\
365 & Artificial intelligence-driven multiomics predictive model for abdominal aortic aneurysm subtypes to identify heterogeneous immune cell infiltration and predict disease progression. &  &  &  & E \\
366 & One-year outcome and durability of pulmonary vein isolation after prospective use of ablation index for catheter ablation in patients with persistent atrial fibrillation. &  &  &  & E \\
367 & Deep learning based MA detection with modified ResNet-50 &  &  &  & E \\
368 & Rationale and design of a multicenter, prospective, diagnostic clinical study: A study protocol for evaluating the diagnostic validation of deep learning-based noninvasive CT-FFR for in-stent restenosis. &  &  &  & E \\
369 & Machine learning-based radiomic features of perivascular adipose tissue in coronary computed tomography angiography predicting inflammation status around atherosclerotic plaque: a retrospective cohort study. &  &  &  & E \\
370 & Machine learning aids clinical decision-making in patients presenting with angina and non-obstructive coronary artery disease. &  &  &  & E \\
371 & Comparison of robot-assisted versus fluoroscopy-assisted minimally invasive transforaminal lumbar interbody fusion for degenerative lumbar spinal diseases: 2-year follow-up. &  &  &  & E \\
372 & Deep Learning for 3D Vascular Segmentation in Phase Contrast Tomography &  &  &  & E \\
373 & Endovascular robotics: technical advances and future directions. &  &  &  & E \\
374 & Improving image quality and in-stent restenosis diagnosis with high-resolution "double-low" coronary CT angiography in patients after percutaneous coronary intervention. &  &  &  & E \\
375 & VWI-APP: Vessel wall imaging-dedicated automated processing pipeline for intracranial atherosclerotic plaque quantification. &  &  &  & E \\
376 & Efficacy of artificial intelligence-based FFR technology for coronary CTA stenosis detection in clinical management of coronary artery disease: a systematic review. &  &  &  & E \\
377 & Radiation-Free Percutaneous Coronary Intervention: Myth or Reality? &  &  &  & E \\
378 & Changes of blood flow in macular zone of patients with diabetic retinopathy at different stages evaluated by optical coherence tomography angiography. &  &  &  & E \\
379 & The Rise of Optical Coherent Tomography in Intracoronary Imaging: An Overview of Current Technology, Limitations, and Future Perspectives. &  &  &  & E \\
380 & Serial analysis of coronary artery disease progression by artificial intelligence assisted coronary computed tomography angiography: early clinical experience. &  &  &  & E \\
381 & Direct head-to-head comparison of convolutional long short-term memory and transformer networks for artificial Intelligence-based quantification of atherosclerotic plaque and stenosis from coronary CT angiography &  &  &  & E \\
382 & One year of atrial fibrillation ablation using pulsed-field pentaspline catheter - learning curves based on a single centre experience &  &  &  & E \\
383 & Coronary atherosclerosis on AI-based plaque analysis in patients with chest pain and calcium score zero. &  &  &  & E \\
384 & Deep learning artery-vein classification in OCT angiography &  &  &  & E \\
385 & Introducing Two Novel Indices for Evaluating Coronary Artery Stenosis: Athero-Inflammatory Index and Athero-Inflammatory Glucose Index. &  &  &  & E \\
386 & Predicting Coronary Stenosis Progression Using Plaque Fatigue From IVUS-Based Thin-Slice Models: A Machine Learning Random Forest Approach. &  &  &  & E \\
387 & Plasma metabolomics disentangles T2DM- and CAD-specific dysmetabolism and identifies potential biomarkers for CAD risk escalation in diabetic patients. &  &  &  & E \\
388 & Optimizing the OCTA layer fusion option for deep learning classification of diabetic retinopathy. &  &  &  & E \\
389 & The role of optical coherence tomography angiography in distinguishing ischemic versus non-ischemic central retinal vein occlusion. &  &  &  & E \\
390 & Prediction of clinical efficacy of acupuncture intervention on upper limb dysfunction after ischemic stroke based on machine learning: a study driven by DSA diagnostic reports data. &  &  &  & E \\
391 & Validation of quantitative perfusion cardiovascular magnetic resonance employing deconvolution techniques with Tofts, modified-Tofts, and Fermi function models against <sup>15</sup>O-water positron emission tomography. &  &  &  & E \\
392 & Machine learning improves prediction of pulmonary thromboembolism and reduces unnecessary computed tomography scans in the emergency department. &  &  &  & E \\
393 & Applications and implementation of generative artificial intelligence in cardiovascular imaging with a focus on ethical and legal considerations: what cardiovascular imagers need to know! &  &  &  & E \\
394 & OCT angiography in the monitoring of vaginal health. &  &  &  & E \\
395 & The diagnostic performance of machine learning-based FFRCT for coronary artery disease: A meta-analysis. &  &  &  & E \\
396 & Glaucoma detection using non-perfused areas in OCTA. &  &  &  & E \\
397 & Evaluation of Bladder Dysfunction Outcomes Among Standardized Bladder Shapes in Children With Spina Bifida. &  &  &  & E \\
398 & Towards a FHIR-Based Data Model for Coronary Angiography Observations. &  &  &  & E \\
399 & Ultrafast Brain MRI with Deep Learning Reconstruction for Suspected Acute Ischemic Stroke. &  &  &  & E \\
400 & A Machine Learning Based Scheme for Dynamic Spectrum Access &  &  &  & E \\
401 & Prediction of disorders with significant coronary lesions using machine learning in patients admitted with chest symptom. &  &  &  & E \\
402 & Deep Learning Using One-stop-shop CT Scan to Predict Hemorrhagic Transformation in Stroke Patients Undergoing Reperfusion Therapy: A Multicenter Study. &  &  &  & E \\
403 & Accurately Identifying Cerebroarterial Stenosis from Angiography Reports Using Natural Language Processing Approaches. &  &  &  & E \\
404 & Abdominal Aortic Thrombus Segmentation in Postoperative Computed Tomography Angiography Images Using Bi-Directional Convolutional Long Short-Term Memory Architecture. &  &  &  & E \\
405 & Can large language models be new supportive tools in coronary computed tomography angiography reporting? &  &  &  & E \\
406 & Deep Learning for Spectrum Prediction in Cognitive Radio Networks: State-of-the-Art, New Opportunities, and Challenges &  &  &  & E \\
407 & Machine Learning-Based CCTA-Guided Intensive Atheroprotective Strategy in a Middle-Aged INOCA Patient With Challenging Arterial Features. &  &  &  & E \\
408 & Consensus for the research and application of artificial intelligence in coronary computed tomography angiography &  &  &  & E \\
409 & Optimizing Coronary Computed Tomography Angiography Using a Novel Deep Learning-Based Algorithm. &  &  &  & E \\
410 & Revolution or risk?-Assessing the potential and challenges of GPT-4V in radiologic image interpretation. &  &  &  & E \\
411 & Artificial intelligence algorithm for detecting myocardial infarction using six-lead electrocardiography. &  &  &  & E \\
412 & Angioarchitectural features of brain arteriovenous malformation presented with seizures. &  &  &  & E \\
413 & Intelligent Transition Mechanisms in 5G DSA: A Logistic Regression Approach to Reuse Factor Adaptation &  &  &  & E \\
414 & Artificial intelligence-based automatic nidus segmentation of cerebral arteriovenous malformation on time-of-flight magnetic resonance angiography. &  &  &  & E \\
415 & Artificial intelligence in coronary computed tomography angiography: Demands and solutions from a clinical perspective. &  &  &  & E \\
416 & New Insights Into Pathogenesis, Diagnosis, and Management of Cardiac Allograft Vasculopathy. &  &  &  & E \\
417 & Clinical Validation of a Machine-Learned, Point-of-Care System to IDENTIFY Functionally Significant Coronary Artery Disease. &  &  &  & E \\
418 & Accuracy of calcified plaque volumes by AI-QCT coronary CT angiography: phantom and non-contrast CT verification. &  &  &  & E \\
419 & Machine Learning Identifies New Predictors on Restenosis Risk after Coronary Artery Stenting in 10,004 Patients with Surveillance Angiography. &  &  &  & E \\
420 & Do plaque-related factors affect the diagnostic performance of an artificial intelligence coronary-assisted diagnosis system? Comparison with invasive coronary angiography. &  &  &  & E \\
421 & Deep-Reinforcement-Learning-Based Distributed Dynamic Spectrum Access in Multiuser Multichannel Cognitive Radio Internet of Things Networks &  &  &  & E \\
422 & Artificial intelligence in aortic CT angiography: current applications and future perspectives. &  &  &  & E \\
423 & Feature-Based Classification of Mild Cognitive Impairment and Alzheimer's Disease Based on Optical Coherence Tomographic Angiographic Image. &  &  &  & E \\
424 & Personalized pulmonary vein antrum isolation guided by left atrial wall thickness for persistent atrial fibrillation. &  &  &  & E \\
425 & ARTIFICIAL INTELLIGENCE-ENHANCED ANALYSIS OF RETINAL VASCULATURE IN AGE-RELATED MACULAR DEGENERATION. &  &  &  & E \\
426 & On-Site Computed Tomography-Derived Fractional Flow Reserve Using a Machine-Learning Algorithm - Clinical Effectiveness in a Retrospective Multicenter Cohort. &  &  &  & E \\
427 & Evolution of Hypodensity on Non-Contrast CT in Correlation with Collaterals in Anterior Circulation Stroke with Successful Endovascular Reperfusion. &  &  &  & E \\
428 & Multi-scale residual denoising GAN model for producing super-resolution CTA images &  &  &  & E \\
429 & Highlights of the society for magnetic resonance angiography 2024 conference. &  &  &  & E \\
430 & Artificial intelligence in coronary computed tomography angiography &  &  &  & E \\
431 & DOVE: Doodled vessel enhancement for photoacoustic angiography super resolution. &  &  &  & E \\
432 & Safety and effectiveness of very-high-power, short-duration ablation in patients with atrial fibrillation: Preliminary results. &  &  &  & E \\
433 & Standardization of Optical Coherence Tomography Angiography Imaging Biomarkers in Diabetic Retinal Disease. &  &  &  & E \\
434 & Radiation dose reduction with deep-learning image reconstruction for coronary computed tomography angiography. &  &  &  & E \\
435 & X23D-Intraoperative 3D Lumbar Spine Shape Reconstruction Based on Sparse Multi-View X-ray Data. &  &  &  & E \\
436 & 30-Day Risk Score for Mortality and Stroke in Patients with Carotid Artery Stenosis Using Artificial Intelligence Based Carotid Plaque Morphology. &  &  &  & E \\
437 & Global trends and prospects of ocular biomarkers in Alzheimer's disease: a bibliometric analysis. &  &  &  & E \\
438 & Robot-Assisted Versus Fluoroscopy-Guided Pedicle Screw Placement in Transforaminal Lumbar Interbody Fusion for Lumbar Degenerative Disease. &  &  &  & E \\
439 & OFF-eNET: An Optimally Fused Fully End-to-End Network for Automatic Dense Volumetric 3D Intracranial Blood Vessels Segmentation &  &  &  & E \\
440 & Prediction of obstructive coronary artery disease in people living with HIV: value of machine learning incorporating HIV-specific factors. &  &  &  & E \\
441 & Deep generative models for vessel segmentation in CT angiography of the brain. &  &  &  & E \\
442 & Machine Learning Efficiency in Predicting Obstructive Coronary Artery Disease in Patients with Non-ST Elevation Acute Coronary Syndrome in the First Hours of Admission. &  &  &  & E \\
443 & Automated assessment of collateral circulation and infarct core: predictors of functional outcomes in acute ischemic stroke following endovascular thrombectomy. &  &  &  & E \\
444 & Diagnostic Performance of Artificial Intelligence-Based Angiography-Derived Non-Hyperemic Pressure Ratio Using Pressure Wire as Reference. &  &  &  & E \\
445 & Anatomical atlas of the upper part of the human head for electroencephalography and bioimpedance applications. &  &  &  & E \\
446 & Quantum-AI Number-Theoretic Frameworks for Post-Quantum Cryptography and Long-Term Information Security &  &  &  & E \\
447 & Automated Deep Learning Analysis for Quality Improvement of CT Pulmonary Angiography. &  &  &  & E \\
448 & Radiomics-Based Precision Phenotyping Identifies Unstable Coronary Plaques From Computed Tomography Angiography. &  &  &  & E \\
449 & Cardiovascular CT and MRI in 2019: Review of Key Articles. &  &  &  & E \\
450 & Artificial intelligence-based automated left ventricular mass quantification from non-contrast cardiac CT scans: correlation with contrast CT and cardiac MRI. &  &  &  & E \\
451 & Enhancing mesenteric vascular imaging with dual-energy CTA: a comparison of DLIR and ASIR-V using low contrast agent and low radiation dose protocols. &  &  &  & E \\
452 & Deep learning image reconstruction and adaptive statistical iterative reconstruction on coronary artery calcium scoring in high risk population for coronary heart disease. &  &  &  & E \\
453 & Fake Product Identification Using Deep Learning &  &  &  & E \\
454 & Ablation index-guided high-power vs. moderate-power cavotricuspid isthmus ablation. &  &  &  & E \\
455 & Metaverse-based deep learning framework for coronary artery stenosis classification using Monte Carlo Dropout-based ResNet-152. &  &  &  & E \\
456 & Radiomics analysis enhances the diagnostic performance of CMR stress perfusion: a proof-of-concept study using the Dan-NICAD dataset. &  &  &  & E \\
457 & Deep Learning Improves the Temporal Reproducibility of Aortic Measurement. &  &  &  & E \\
458 & Multi-Energy Evaluation of Image Quality in Spectral CT Pulmonary Angiography Using Different Strength Deep Learning Spectral Reconstructions. &  &  &  & E \\
459 & Time- and Frequency-Domain Dynamic Spectrum Access: Learning Cyclic Medium Access Patterns in Partially Observable Environments &  &  &  & E \\
460 & Five-year systemic complications in diabetic retinopathy with integrated optical coherence tomography angiography and glycated hemoglobin. &  &  &  & E \\
461 & DeepCA: Deep Learning-Based 3D Coronary Artery Tree Reconstruction from Two 2D Non-Simultaneous X-Ray Angiography Projections &  &  &  & E \\
462 & Takotsubo syndrome in a patient with adrenal insufficiency, severe hyponatremia and coexistent coronary artery disease. &  &  &  & E \\
463 & Abstract 381: Artificial Intelligence as a Diagnostic Aid: A Systematic Review and Meta‐Analysis of Human‐AI Collaboration for Intracranial Aneurysm Detection &  &  &  & E \\
464 & Primary User Activity Prediction in DSA Networks using Recurrent Structures &  &  &  & E \\
465 & Intracoronary Optical Coherence Tomography: Technological Innovations and Clinical Implications in Cardiology. &  &  &  & E \\
466 & Is the Acetabular Cup Orientation Different in Robot-Assisted and Conventional Total Hip Arthroplasty With Right-Handed Surgeons Using an Anterolateral Approach? &  &  &  & E \\
467 & Chest pain and AI interpretation: The essential role of clinical judgment and quality data. &  &  &  & E \\
468 & Quantitative plaque analysis with A.I.-augmented CCTA in end-stage renal disease and complex CAD. &  &  &  & E \\
469 & ChatGPT vs Gemini: Comparative Accuracy and Efficiency in CAD-RADS Score Assignment from Radiology Reports. &  &  &  & E \\
470 & Precise Staging of Diabetic Retinopathy Through Machine Learning Analysis of Leakage Source Characteristics: A Radiomics-Driven Approach Integrating OCT-A and FFA. &  &  &  & E \\
471 & Advances in Diagnosis, Therapy, and Prognosis of Coronary Artery Disease Powered by Deep Learning Algorithms. &  &  &  & E \\
472 & Spectral Temporal Graph Neural Network for Multivariate Time-series Forecasting &  &  &  & E \\
473 & Radiomic Assessment of Epicardial Adipose Tissue for the Prediction of Non-Calcified Coronary Atherosclerotic Plaques. &  &  &  & E \\
474 & Feasibility Study of Triple-low CCTA for Coronary Artery Disease Screening Combining Contrast Enhancement Boost and Deep Learning Reconstruction. &  &  &  & E \\
475 & Advances in retinopathy of prematurity imaging. &  &  &  & E \\
476 & The Characteristic of Optical Coherence Tomography Angiography and Retinal Arteries Angle in Familial Exudative Vitreoretinopathy with Inner Retinal Layer Persistence. &  &  &  & E \\
477 & Reduced macular thickness and vascular density in abnormal glucose metabolism patients: A meta-analysis of optical coherence tomography (OCT) and OCT angiography studies. &  &  &  & E \\
478 & The Impact of a Neuroradiologist on the Report of a Real-World CT Perfusion Imaging Map Derived from Artificial Intelligence/Machine Learning-Driven Software. &  &  &  & E \\
479 & Functional Acetabular Component Positioning During Direct Anterior Approach Hip Arthroplasty Using a Novel Three-Dimensional Virtual Mesh Imaging System With Fluoroscopy. &  &  &  & E \\
480 & Determinants of Progression and Regression of Subclinical Atherosclerosis: A Temporal Perspective. &  &  &  & E \\
481 & Effect of surgical intraocular pressure lowering on retinal structures - nerve fibre layer, foveal avascular zone, peripapillary and macular vessel density: 1 year results. &  &  &  & E \\
482 & Bifascicular Block Associated With Myocardial Infarction: A Marker of Proximal Left Anterior Descending Artery Occlusion Confirmed by the Artificial Intelligence-Based Smartphone App Queen of Hearts. &  &  &  & E \\
483 & Development \& Validation of a Machine Learning Model That Uses Voice to Predict Aspiration Risk. &  &  &  & E \\
484 & Impact of tracer uptake rate on quantification accuracy of myocardial blood flow in PET: A simulation study. &  &  &  & E \\
485 & Weakly-supervised convolutional neural networks of renal tumor segmentation in abdominal CTA images. &  &  &  & E \\
486 & Automated detection of large vessel occlusion using deep learning: a pivotal multicenter study and reader performance study. &  &  &  & E \\
487 & Diagnostic Accuracy In Coronary CT Angiography Analysis: Artificial Intelligence Versus Visual Assessment &  &  &  & E \\
488 & Association of Subregional Quantitative Ultra-widefield Fluorescence Angiography Characteristics With the Occurrence of Diabetic Macular Edema and Proliferative Diabetic Retinopathy. &  &  &  & E \\
489 & Validation of angiography-derived Murray law-based quantitative flow reserve (μQFR) against pressure-derived instantaneous wave-free ratio for assessing coronary lesions, a single-center study in Egypt. &  &  &  & E \\
490 & Optical coherence tomography angiography in cardiovascular disease. &  &  &  & E \\
491 & FSTIF-UNet: A Deep Learning-Based Method Towards Automatic Segmentation of Intracranial Aneurysms in Un-Reconstructed 3D-RA. &  &  &  & E \\
492 & Calcium Scoring at Coronary CT Angiography Using Deep Learning. &  &  &  & E \\
493 & Cerebral Angiography under Artificial Intelligence Algorithm in the Design of Nursing Cooperation Plan for Intracranial Aneurysm Patients in Craniotomy Clipping. &  &  &  & E \\
494 & Machine learning and statistical approaches for classification of risk of coronary artery disease using plasma cytokines. &  &  &  & E \\
495 & Machine Learning From Quantitative Coronary Computed Tomography Angiography Predicts Ischemia And Impaired Myocardial Blood Flow &  &  &  & E \\
496 & Machine learning-based coronary heart disease diagnosis model for type 2 diabetes patients. &  &  &  & E \\
497 & Quantifying the calcification of abdominal aorta and major side branches with deep learning. &  &  &  & E \\
498 & Radiomics of pericoronary adipose tissue on computed tomography angiography predicts coronary heart disease in patients with type 2 diabetes mellitus. &  &  &  & E \\
499 & Anatomy-Informed Deep Learning for Abdominal Aortic Aneurysm Segmentation &  &  &  & E \\
500 & Detection of intracranial aneurysms using deep learning-based CAD system: usefulness of the scores of CNN's final layer for distinguishing between aneurysm and infundibular dilatation. &  &  &  & E \\
501 & Collaborative Multiscale and Wavelet-Based Fusion Network for Leakage Area Semantic Segmentation of Ultrawide Field Fluorescein Angiography. &  &  &  & E \\
502 & A retrieval-augmented chatbot based on GPT-4 provides appropriate differential diagnosis in gastrointestinal radiology: a proof of concept study. &  &  &  & E \\
503 & Discovering hidden patterns: Association rules for cardiovascular diseases in type 2 diabetes mellitus. &  &  &  & E \\
504 & TREAT-Netv2: regional wall motion-informed video-tabular fusion for ACS treatment prediction. &  &  &  & E \\
505 & Artificial intelligence of stenosis on coronary CTA: real world comparison with quantitative coronary angiography. &  &  &  & E \\
506 & Artificial Intelligence Iterative Reconstruction Algorithm Combined with Low-Dose Aortic CTA for Preoperative Access Assessment of Transcatheter Aortic Valve Implantation: A Prospective Cohort Study. &  &  &  & E \\
507 & Taxonomy of Acute Stroke: Imaging, Processing, and Treatment. &  &  &  & E \\
508 & Editorial: Applications of optical coherence tomography angiography (OCTA) in ocular complications of diabetes mellitus. &  &  &  & E \\
509 & Correlation between the high-density lipoprotein-to-low-density lipoprotein ratio and contrast-induced acute kidney injury: an analysis based on the in-hospital AKI electronic monitoring system. &  &  &  & E \\
510 & Improving the diagnostic performance of computed tomography angiography for intracranial large arterial stenosis by a novel super-resolution algorithm based on multi-scale residual denoising generative adversarial network. &  &  &  & E \\
511 & Application of an interpretable machine learning model based on optimal feature selection for predicting triple-vessel coronary disease: a multicenter retrospective study. &  &  &  & E \\
512 & Diagnostic performance of machine-learning-based computed fractional flow reserve (FFR) derived from coronary computed tomography angiography for the assessment of myocardial ischemia verified by invasive FFR. &  &  &  & E \\
513 & Current Status and Future Perspectives of Optic Nerve Imaging in Glaucoma. &  &  &  & E \\
514 & Deep learning-based head and neck segmentation in CT angiography &  &  &  & E \\
515 & Machine learning of clinical variables and coronary artery calcium scoring for the prediction of obstructive coronary artery disease on coronary computed tomography angiography: analysis from the CONFIRM registry. &  &  &  & E \\
516 & Quantification of Aortic Valve Fibrotic and Calcific Tissue from CTA: Prospective Comparison with Histology. &  &  &  & E \\
517 & Dynamic simulation of knee joint mechanics: individualized multi-moment finite element modelling of patellar tendon stress during landing. &  &  &  & E \\
518 & Sequential vessel segmentation via deep channel attention network. &  &  &  & E \\
519 & Artificial intelligence-based quantitative coronary angiography of major vessels using deep-learning. &  &  &  & E \\
520 & Non - invasive modelling methodology for the diagnosis of coronary artery disease using fuzzy cognitive maps. &  &  &  & E \\
521 & Research on A Chiplet-based DSA (Domain-Specific Architectures) Scalable Convolutional Acceleration Architecture &  &  &  & E \\
522 & Classification of Cardiovascular Diseases through Retinal Fundus Images using Deep Learning &  &  &  & E \\
523 & Accelerating non-contrast MR angiography of the thoracic aorta using compressed SENSE with deep learning reconstruction. &  &  &  & E \\
524 & Vessel diameters of 14 basal cerebral arteries assessed in 1000 digital subtraction angiographies. &  &  &  & E \\
525 & Segregation of neuronal-vascular components in a retinal nerve fiber layer for thickness measurement using OCT and OCT angiography. &  &  &  & E \\
526 & Incorporating clinical parameters to improve the accuracy of angiography-derived computed fractional flow reserve<sup></sup>. &  &  &  & E \\
527 & Fully automated guideline-compliant diameter measurements of the thoracic aorta on ECG-gated CT angiography using deep learning. &  &  &  & E \\
528 & Accelerated Intravascular Ultrasound Imaging using Deep Reinforcement Learning &  &  &  & E \\
529 & Differentiation of Diabetic Status Using Statistical and Machine Learning Techniques on Optical Coherence Tomography Angiography Images. &  &  &  & E \\
530 & ASL 4D MRA Intracranial Vessel Segmentation With Deep Learning U-Nets. &  &  &  & E \\
531 & A multi-backbone cascade and morphology-aware segmentation network for complex morphological X-ray coronary artery images. &  &  &  & E \\
532 & Establishing a Classification System for Predicting Flow-Limiting Dissection After Balloon Angioplasty Using Explainable Machine-Learning Models: A Multicenter Retrospective Cohort Study. &  &  &  & E \\
533 & Noninvasive Machine Learning Screening Model for Dacryocystitis Based on Ocular Surface Indicators. &  &  &  & E \\
534 & Deep learning reconstruction allows for usage of contrast agent of lower concentration for coronary CTA than filtered back projection and hybrid iterative reconstruction. &  &  &  & E \\
535 & LUMEN-A deep learning pipeline for analysis of the 3D morphology of the cerebral lenticulostriate arteries from time-of-flight 7T MRI. &  &  &  & E \\
536 & Translating color fundus photography to indocyanine green angiography using deep-learning for age-related macular degeneration screening. &  &  &  & E \\
537 & PET Myocardial Flow Reserve Estimation from 4D-Coronary-CT using Deep Neural Network. &  &  &  & E \\
538 & Automatic Detection and Classification of Diabetic Retinopathy from Optical Coherence Tomography Angiography Images using Deep Learning - A Review &  &  &  & E \\
539 & CT Coronary Angiography Fractional Flow Reserve: New Advances in the Diagnosis and Treatment of Coronary Artery Disease. &  &  &  & E \\
540 & Comparison of Cryoballoon and Ablation Index-Guided Radiofrequency Ablation in Paroxysmal Atrial Fibrillation. &  &  &  & E \\
541 & Optimizing deep learning for brain vessel segmentation in CT angiography: a study on domain adaptation and generalizability &  &  &  & E \\
542 & Prediction of coronary artery calcium scoring from surface electrocardiogram in atherosclerotic cardiovascular disease: a pilot study. &  &  &  & E \\
543 & Advanced Preoperative Planning Techniques in the Management of Complex Proximal Humerus Fractures. &  &  &  & E \\
544 & Contemporary, non-invasive imaging diagnosis of chronic coronary artery disease. &  &  &  & E \\
545 & Segmentation and Classification of Heart Angiographic Images Using Machine Learning Techniques. &  &  &  & E \\
546 & Impact of the ADMIRE reconstruction algorithm combined with the Sa36 kernel on quantitative measurement of coronary artery calcification in AI: a single-arm prospective study. &  &  &  & E \\
547 & Retinal microvascular alterations in macular telangiectasia type 2: An optical coherence tomography angiography study. &  &  &  & E \\
548 & Region guided mask R-CNN with Haralick ResNet fusion for accurate coronary artery disease detection in computed tomography angiography images. &  &  &  & E \\
549 & Enhancing semi-supervised learning for fine-grained 3D cerebrovascular segmentation with cross-consistency and uncertainty estimation. &  &  &  & E \\
550 & Assessment of atherosclerotic plaque burden: comparison of AI-QCT versus SIS, CAC, visual and CAD-RADS stenosis categories. &  &  &  & E \\
551 & Peri-procedural anesthesia and patient pain experience in pulmonary vein isolation by means of very high-power short-duration radiofrequency ablation. &  &  &  & E \\
552 & Dynamic spectrum access based on double deep Q-network and convolution neural network &  &  &  & E \\
553 & Improving Image Resolution of Whole-Heart Coronary MRA Using Convolutional Neural Network. &  &  &  & E \\
554 & Validation of deep-learning image reconstruction for coronary computed tomography angiography: Impact on noise, image quality and diagnostic accuracy. &  &  &  & E \\
555 & Should they stay or should they go? Stroke transfers across a hospital network pre- and post-implementation of an automated image interpretation and communication platform. &  &  &  & E \\
556 & Coronary CT Fractional Flow Reserve before Transcatheter Aortic Valve Replacement: Clinical Outcomes. &  &  &  & E \\
557 & Noninvasive Synthesis of Multiframe Ultra-Widefield Fluorescein Angiography from Color Fundus Photographs. &  &  &  & E \\
558 & Prediction of short-term adverse clinical outcomes of acute pulmonary embolism using conventional machine learning and deep Learning based on CTPA images. &  &  &  & E \\
559 & Use of a deep-learning-based lumen extraction method to detect significant stenosis on coronary computed tomography angiography in patients with severe coronary calcification. &  &  &  & E \\
560 & A fusion model of deep learning and conventional features based on computed tomography angiography of carotid plaque for predicting the risk of acute ischemic stroke. &  &  &  & E \\
561 & Artificial Intelligence and Cardiovascular Genetics. &  &  &  & E \\
562 & Enhancing Early Detection of Brain Aneurysms: A CNN-Driven, Real-Time Approach with Angiography Imaging &  &  &  & E \\
563 & Weisfeiler and Leman go Machine Learning: The Story so far &  &  &  & E \\
564 & Impact of wrist-ankle acupuncture on propofol dosage under the dual monitoring of density spectrum array and anesthesia consciousness index in elderly patients undergoing urologic surgery: a sham-controlled randomized clinical trial. &  &  &  & E \\
565 & Artificial Intelligence Based Multimodality Imaging: A New Frontier in Coronary Artery Disease Management. &  &  &  & E \\
566 & Diagnostic accuracy of 3D deep-learning-based fully automated estimation of patient-level minimum fractional flow reserve from coronary computed tomography angiography. &  &  &  & E \\
567 & Current Trends and Perspectives of Pressure Wire-Based Coronary Artery Bypass Grafting. &  &  &  & E \\
568 & Diagnostic Performance Of A Novel Onsite Deep Learning-based Coronary CT Angiography-derived Fractional Flow Reserve For Detecting Functionally Significant Coronary Stenosis: A Comparison With Instantaneous Wave-Free Ratio &  &  &  & E \\
569 & Using Neural Networks Algorithm in Ischemic Stroke Diagnosis: A Systematic Review. &  &  &  & E \\
570 & ChatFFA: An ophthalmic chat system for unified vision-language understanding and question answering for fundus fluorescein angiography. &  &  &  & E \\
571 & Generative Artificial Intelligence to Automate Cerebral Perfusion Mapping in Acute Ischemic Stroke from Non-contrast Head Computed Tomography Images: Pilot Study. &  &  &  & E \\
572 & Identification of Blood Vessels from Anigography Images Using CNN &  &  &  & E \\
573 & Machine learning-enhanced echocardiography for screening coronary artery disease. &  &  &  & E \\
574 & Comparison of DLIR and ASIR-V algorithms for virtual monoenergetic imaging in carotid CTA under a triple-low protocol. &  &  &  & E \\
575 & Usefulness of metal artifact reduction on CT angiography after massive coil embolization in peripheral AVM. &  &  &  & E \\
576 & Accuracy of an Artificial Intelligence Deep Learning Algorithm Implementing a Recurrent Neural Network With Long Short-term Memory for the Automated Detection of Calcified Plaques From Coronary Computed Tomography Angiography. &  &  &  & E \\
577 & Prognostic value of CT myocardial perfusion imaging and CT-derived fractional flow reserve for major adverse cardiac events in patients with coronary artery disease. &  &  &  & E \\
578 & Optical coherence tomography (OCT) and OCT angiography: Technological development and applications in brain science. &  &  &  & E \\
579 & Phenotypic clusters of patients with Takayasu arteritis predict disease activity at presentation, damage and future vascular complications. &  &  &  & E \\
580 & Defense against PUE Attacks in DSA Networks Using GAN Based Learning &  &  &  & E \\
581 & The Role of Non-invasive Tests in Pulmonary Embolism. &  &  &  & E \\
582 & Enhanced clinical photoacoustic vascular imaging through a skin localization network and adaptive weighting &  &  &  & E \\
583 & Application of Artificial Intelligence in Vulnerable Carotid Atherosclerotic Plaque Assessment-A Scoping Review. &  &  &  & E \\
584 & Vessel-specific angiography-derived index of microcirculatory resistance in an all-comer population undergoing percutaneous coronary intervention, a PIONEER IV trial substudy. &  &  &  & E \\
585 & A systematic review of artificial intelligence tools for chronic pulmonary embolism on CT pulmonary angiography. &  &  &  & E \\
586 & Diagnostic accuracy of automated occlusion detection in CT angiography using e-CTA. &  &  &  & E \\
587 & Addressing Limited Generalizability in Artificial Intelligence-Based Brain Aneurysm Detection for Computed Tomography Angiography: Development of an Externally Validated Artificial Intelligence Screening Platform. &  &  &  & E \\
588 & ERS International Congress 2023: highlights from the Clinical Techniques, Imaging and Endoscopy Assembly. &  &  &  & E \\
589 & Artificial intelligence techniques for cardiovascular disease diagnosis via X-ray sensor-based coronary angiography: A bibliometric and systematic review. &  &  &  & E \\
590 & Circulating organokines in coronary artery disease and metabolic syndrome: FABP4, adiponectin, irisin, FSTL1. &  &  &  & E \\
591 & Multidisciplinary In-Depth Investigation in a Young Athlete Suffering from Syncope Caused by Myocardial Bridge. &  &  &  & E \\
592 & Classification models for assessing coronary artery disease instances using clinical and biometric data: an explainable man-in-the-loop approach. &  &  &  & E \\
593 & An artificial intelligence algorithm for pulmonary embolism detection on polychromatic computed tomography: performance on virtual monochromatic images. &  &  &  & E \\
594 & Exploring the predictive values of CRP and lymphocytes in coronary artery disease based on a machine learning and Mendelian randomization. &  &  &  & E \\
595 & High-Speed On-Site Deep Learning-Based FFR-CT Algorithm: Evaluation Using Invasive Angiography as the Reference Standard. &  &  &  & E \\
596 & Artificial Intelligence for Cardiovascular Care-Part 1: Advances: JACC Review Topic of the Week. &  &  &  & E \\
597 & Estimating degree of coronary artery stenosis severity from non-contrast CT images using CNN &  &  &  & E \\
598 & Non-invasive imaging as the cornerstone of cardiovascular precision medicine. &  &  &  & E \\
599 & Development and validation of an AMR-based predictive model for post-PCI upper gastrointestinal bleeding in NSTEMI patients. &  &  &  & E \\
600 & Value Proposition of FDA-Approved Artificial Intelligence Algorithms for Neuroimaging. &  &  &  & E \\
601 & Review of Vessel Segmentation and Stenosis classification in X-ray Coronary Angiography &  &  &  & E \\
602 & Radiologist Worklist Reprioritization Using Artificial Intelligence: Impact on Report Turnaround Times for CTPA Examinations Positive for Acute Pulmonary Embolism. &  &  &  & E \\
603 & Comparison of image quality in 40 keV virtual monoenergetic images of dual-energy CT pulmonary angiography using deep learning and iterative reconstruction algorithms under optimized low dose scanning protocols. &  &  &  & E \\
604 & Performance of image-based deep learning models for aortic dissection segmentation and diagnosis: a systematic review and meta-analysis. &  &  &  & E \\
605 & Multi-step deep neural network for identifying subfascial vessels in a dorsal skinfold window chamber model. &  &  &  & E \\
606 & AI-Enhanced Fluorescein Angiography Detection of Diabetes-Induced Silent Retinal Capillary Dropout and RNA-Seq Identification of Pre-Symptomatic Biomarkers. &  &  &  & E \\
607 & A lung CT vision foundation model facilitating disease diagnosis and medical imaging. &  &  &  & E \\
608 & Improved visualization of intracranial distal arteries with multiple 2D slice dynamic ASL-MRA and super-resolution convolutional neural network. &  &  &  & E \\
609 & A machine learning pipeline for internal anatomical landmark embedding based on a patient surface model. &  &  &  & E \\
610 & Advances in Imaging and Diagnosis of Abdominal Aortic Aneurysm: A Shift in Clinical Practice. &  &  &  & E \\
611 & A Deep-Learning-Based Framework for Automatic Segmentation and Labelling of Intracranial Artery &  &  &  & E \\
612 & Modern Imaging Techniques for Percutaneous Coronary Intervention Guidance: A Focus on Intravascular Ultrasound and Optical Coherence Tomography. &  &  &  & E \\
613 & Door to Puncture in Large Vessel Occlusions Pre- and Postimplementation of an Automated Image Interpretation and Communication Platform: A Single Center Study. &  &  &  & E \\
614 & Right-to-left ventricle ratio determined by machine learning algorithms on CT pulmonary angiography images predicts prolonged ICU length of stay in operated chronic thromboembolic pulmonary hypertension. &  &  &  & E \\
615 & Combining generative modelling and semi-supervised domain adaptation for whole heart cardiovascular magnetic resonance angiography segmentation. &  &  &  & E \\
616 & Implementing an AI algorithm in the clinical setting: a case study for the accuracy paradox. &  &  &  & E \\
617 & Real-World Experience with Artificial Intelligence-Based Triage in Transferred Large Vessel Occlusion Stroke Patients. &  &  &  & E \\
618 & Deep learning for artery-vein classification in optical coherence tomography angiography. &  &  &  & E \\
619 & Deep learning accurately and reliably segments pelvic vascular structure in CT scans of gynecologic cancer patients. &  &  &  & E \\
620 & Perivascular adipose tissue as a source of therapeutic targets and clinical biomarkers. &  &  &  & E \\
621 & Artificial Intelligence and Machine Learning in Reconstructive Microsurgery. &  &  &  & E \\
622 & Feasibility of U-Net model for cerebral arteries segmentation with low-dose computed tomography angiographic images with pre-processing methods. &  &  &  & E \\
623 & Cognition-Eye-Brain Connection in Alzheimer's Disease Spectrum Revealed by Multimodal Imaging. &  &  &  & E \\
624 & Failure of standard contemporary ST-elevation myocardial infarction electrocardiogram criteria to reliably identify acute occlusion of the left anterior descending coronary artery. &  &  &  & E \\
625 & Comparison of the methods of operation of the artificial intelligence system in the ultra-high sensitivity mode for the autonomous description of chest X-rays without pathology &  &  &  & E \\
626 & PMMNet: A Dual Branch Fusion Network of Point Cloud and Multi-View for Intracranial Aneurysm Classification and Segmentation &  &  &  & E \\
627 & Predictors of discordance between CT-derived fractional flow reserve (CT-FFR) and △CT-FFR in deep coronary myocardial bridging. &  &  &  & E \\
628 & Probability-Based Two-Stage CNN for Pulmonary Embolism Detection in Computed Tomography Pulmonary Angiography &  &  &  & E \\
629 & Exploring single-mode VCSEL wavelength tuning for low-cost 3D optical coherence tomography and OCT angiography. &  &  &  & E \\
630 & Machine learning applications for vascular stenosis detection in computed tomography angiography: a systematic review and meta-analysis &  &  &  & E \\
631 & Artificial intelligence machine learning-based coronary CT fractional flow reserve (CT-FFR<sub>ML</sub>): Impact of iterative and filtered back projection reconstruction techniques. &  &  &  & E \\
632 & Artificial Intelligence-Based Prediction of Contrast Medium Doses for Computed Tomography Angiography Using Optimized Clinical Parameter Sets. &  &  &  & E \\
633 & Prediction model for the risk of vitreous haemorrhage after vitrectomy combined with intraocular injection for the treatment of proliferative diabetic retinopathy. &  &  &  & E \\
634 & MultiD4CAD: Multimodal Dataset composed of CT and Clinical Features for Coronary Artery Disease Analysis. &  &  &  & E \\
635 & Deep Learning Image Analysis of Macular Optical Coherence Tomography Angiography Images for Detection of Progression in Glaucoma &  &  &  & E \\
636 & Multicenter evaluation of interpretable AI for coronary artery disease diagnosis from PET biomarkers. &  &  &  & E \\
637 & <i>Detection of Large Vessel Occlusion Using AI: Evaluating Performance of RapidAI LVO vs Viz.ai LVO in 1,589 Consecutive Code Strokes</i>" (DUEL). &  &  &  & E \\
638 & Diabetes, Atherosclerosis, and Stenosis by AI. &  &  &  & E \\
639 & Recent advances and applications of optical coherence tomography angiography in diabetic retinopathy. &  &  &  & E \\
640 & CT Angiography Radiomics Combining Traditional Risk Factors to Predict Brain Arteriovenous Malformation Rupture: a Machine Learning, Multicenter Study. &  &  &  & E \\
641 & Instability of Variable-selection Algorithms Used to Identify True Predictors of an Outcome in Intermediate-dimension Epidemiologic Studies. &  &  &  & E \\
642 & Improving cerebral microvascular image quality of optical coherence tomography angiography with deep learning-based segmentation. &  &  &  & E \\
643 & The impact of lipoprotein(a) on coronary atherosclerotic plaque phenotype in primary prevention. &  &  &  & E \\
644 & High-Risk Carotid Lesion Segmentation: Advancing Stroke Risk Detection With Deep Learning. &  &  &  & E \\
645 & Global Hotspots and Research Frontiers of Moyamoya Disease and Radiological Features: A Bibliometric Analysis. &  &  &  & E \\
646 & Application of 3T contrast-enhanced Dixon water-fat separation coronary magnetic resonance angiography with artificial intelligence-assisted compressed sensing on coronary artery aneurysm in children with Kawasaki disease. &  &  &  & E \\
647 & Semi-supervised region-connectivity-based cerebrovascular segmentation for time-of-flight magnetic resonance angiography image. &  &  &  & E \\
648 & Artificial intelligence-driven quantitative analysis of CT morphological differences between chronic thromboembolic pulmonary hypertension and chronic thromboembolic disease. &  &  &  & E \\
649 & Deep Shape Features for Predicting Future Intracranial Aneurysm Growth. &  &  &  & E \\
650 & Interpretable artificial intelligence model for predicting heart failure severity after acute myocardial infarction. &  &  &  & E \\
651 & Health Plan-Based Atherosclerosis Imaging for Members with Diabetics: Utilization and Impact on Subsequent Episodes of Care. &  &  &  & E \\
652 & Machine learning-based identification of symptomatic carotid atherosclerotic plaques with dual-energy computed tomography angiography. &  &  &  & E \\
653 & Application of Deep Learning for Automated Detection of Polypoidal Choroidal Vasculopathy in Spectral Domain Optical Coherence Tomography. &  &  &  & E \\
654 & Construction and Interpretability of a Multimodal Deep Learning Model of Electronystagmography-Optical Coherence Tomography Angiography for Early Screening of Alzheimer's Disease. &  &  &  & E \\
655 & Author Correction: A deep learning model for identifying diabetic retinopathy using optical coherence tomography angiography. &  &  &  & E \\
656 & Prediction of Post-Treatment Visual Acuity in Age-Related Macular Degeneration Patients With an Interpretable Machine Learning Method. &  &  &  & E \\
657 & Strategies to combine 3D vasculature and brain CTA with deep neural networks: Application to LVO. &  &  &  & E \\
658 & From Vision to Illumination: The Promethean Journey of Optical Coherence Tomography in Cardiology. &  &  &  & E \\
659 & JOINT MOTION CORRECTION AND 3D SEGMENTATION WITH GRAPH-ASSISTED NEURAL NETWORKS FOR RETINAL OCT. &  &  &  & E \\
660 & A Robust Blood Vessel Segmentation Technique for Angiographic Images Employing Multi-Scale Filtering Approach. &  &  &  & E \\
661 & Artificial intelligence-based prediction of diabetic retinopathy evolution (EviRed): protocol for a prospective cohort. &  &  &  & E \\
662 & Predictive performance of machine learning models for kidney complications following coronary interventions: a systematic review and meta-analysis. &  &  &  & E \\
663 & Image to Image Translation for DSA Image Augmentation on Multimodal Vascular Image Alignment &  &  &  & E \\
664 & Cortical Vessel Segmentation for Neuronavigation Using Vesselness-Enforced Deep Neural Networks &  &  &  & E \\
665 & Illuminating the Invisible: Translational Molecular Imaging in Cardiovascular Medicine. &  &  &  & E \\
666 & DSA-DAMB: A Dynamic Sparse Attention and Defect-Aware Memory Approach for Rare Industrial Defects &  &  &  & E \\
667 & Computerized Differentiation of Growth Status for Abdominal Aortic Aneurysms: A Feasibility Study. &  &  &  & E \\
668 & An Effective Skin Disease Segmentation Model Based on Deep Convolutional Neural Network &  &  &  & E \\
669 & Total Knee Arthroplasty Kinematics Predict Patient-Reported Outcome Measures: Implications for Clinical Kinematic Examinations. &  &  &  & E \\
670 & Deep learning-assistance significantly increases the detection sensitivity of neurosurgery residents for intracranial aneurysms in subarachnoid hemorrhage. &  &  &  & E \\
671 & Cross-modality image synthesis from TOF-MRA to CTA using diffusion-based models. &  &  &  & E \\
672 & Stretched reconstruction based on 2D freehand ultrasound for peripheral artery imaging. &  &  &  & E \\
673 & Digital twins and simulations in transcatheter coronary and structural heart interventions. &  &  &  & E \\
\end{longtable}
